(1) 热容 \(C>0\) 的均温反应器置于温度 \(T_a>0\) 的巨大热库，反应热释放率 \(Q_0\exp[-E/(RT)]\)，壁面热损率 \(H(T-T_a)\)，其中所有常数正。建立温度能量方程；在 \(|T-T_a|\ll T_a\) 的指数线性化下定义 \(\theta=E(T-T_a)/(RT_a^2)\)、\(\tau=Ht/C\)，将方程化为 \(\theta'=\delta e^\theta-\theta\)，明确 \(\delta\)。
(2) 求全部非负定常温度的存在数目、Lambert \(W\) 闭式和临界参数；当 \(\delta=1/e\) 时求临界温升、证明双重根并从两个分支的局部展开得到临界指数。
(3) 对 \(0<\delta<1/e\) 线性化讨论低温和高温定常支的稳定性、扰动特征时间与物理点火阈值，并在 \(\delta>1/e\) 的理想指数模型中证明解会在有限无量纲时间内发散。
(4) 给出一维Lyapunov函数和耗散恒等式；对比原始未线性化的Arrhenius热释放，证明原物理模型不会达到无限温度，并解释指数近似的有效域。
